\chapter*{Алфавиты в математике}\addcontentsline{toc}{chapter}{Алфавиты в математике}\markright{Алфавиты}

Буква сама по себе не задаёт объект: её смысл устанавливается в определении или условии. Например, \(i\) может быть индексом, а в главе о комплексных числах — мнимой единицей. Прописные и строчные буквы различаются.

\begingroup\small\linespread{1}\selectfont\renewcommand{\arraystretch}{1.0}\arrayrulecolor{Teal}

\section*{Латинский алфавит}

Названия приведены по английскому алфавиту. В русских лекциях математические буквы часто читаются иначе: например, «а», «бэ», «цэ», «икс», «игрек». Это не меняет их роли в формуле.

\begin{longtable}{|p{.23\linewidth}|p{.23\linewidth}|p{.44\linewidth}|}\hline Прописная&Строчная&Название\\\hline\endfirsthead\hline Прописная&Строчная&Название\\\hline\endhead

\(A\)&\(a\)&эй\\\hline

\(B\)&\(b\)&би\\\hline

\(C\)&\(c\)&си\\\hline

\(D\)&\(d\)&ди\\\hline

\(E\)&\(e\)&и\\\hline

\(F\)&\(f\)&эф\\\hline

\(G\)&\(g\)&джи\\\hline

\(H\)&\(h\)&эйч\\\hline

\(I\)&\(i\)&ай\\\hline

\(J\)&\(j\)&джей\\\hline

\(K\)&\(k\)&кей\\\hline

\(L\)&\(l\)&эл\\\hline

\(M\)&\(m\)&эм\\\hline

\(N\)&\(n\)&эн\\\hline

\(O\)&\(o\)&оу\\\hline

\(P\)&\(p\)&пи\\\hline

\(Q\)&\(q\)&кью\\\hline

\(R\)&\(r\)&ар\\\hline

\(S\)&\(s\)&эс\\\hline

\(T\)&\(t\)&ти\\\hline

\(U\)&\(u\)&ю\\\hline

\(V\)&\(v\)&ви\\\hline

\(W\)&\(w\)&дабл-ю\\\hline

\(X\)&\(x\)&экс\\\hline

\(Y\)&\(y\)&уай\\\hline

\(Z\)&\(z\)&зед\\\hline

\end{longtable}

\clearpage\section*{Греческий алфавит}

Варианты начертания одной буквы: \(\epsilon,\varepsilon\); \(\phi,\varphi\); \(\theta,\vartheta\); \(\rho,\varrho\); \(\pi,\varpi\); \(\sigma,\varsigma\). Смысл зависит от контекста; разные варианты могут обозначать разные объекты, если автор так договорился.

\begin{longtable}{|p{.23\linewidth}|p{.23\linewidth}|p{.44\linewidth}|}\hline Прописная&Строчная&Название\\\hline\endfirsthead\hline Прописная&Строчная&Название\\\hline\endhead

\(A\)&\(\alpha\)&альфа\\\hline

\(B\)&\(\beta\)&бета\\\hline

\(\Gamma\)&\(\gamma\)&гамма\\\hline

\(\Delta\)&\(\delta\)&дельта\\\hline

\(E\)&\(\varepsilon\)&эпсилон\\\hline

\(Z\)&\(\zeta\)&дзета\\\hline

\(H\)&\(\eta\)&эта\\\hline

\(\Theta\)&\(\theta\)&тета\\\hline

\(I\)&\(\iota\)&йота\\\hline

\(K\)&\(\kappa\)&каппа\\\hline

\(\Lambda\)&\(\lambda\)&лямбда\\\hline

\(M\)&\(\mu\)&мю\\\hline

\(N\)&\(\nu\)&ню\\\hline

\(\Xi\)&\(\xi\)&кси\\\hline

\(O\)&\(o\)&омикрон\\\hline

\(\Pi\)&\(\pi\)&пи\\\hline

\(P\)&\(\rho\)&ро\\\hline

\(\Sigma\)&\(\sigma\)&сигма\\\hline

\(T\)&\(\tau\)&тау\\\hline

\(\Upsilon\)&\(\upsilon\)&ипсилон\\\hline

\(\Phi\)&\(\varphi\)&фи\\\hline

\(X\)&\(\chi\)&хи\\\hline

\(\Psi\)&\(\psi\)&пси\\\hline

\(\Omega\)&\(\omega\)&омега\\\hline

\end{longtable}

\endgroup
